# Exponential and Erlang Kill Processes for Evidence Accumulation Models

## Objective

The aim is to augment a standard evidence accumulation model with a stochastic timeout process that can produce omissions, or time-based guesses, without treating these outcomes as unrelated contaminants. The response process remains interpretable in terms of conventional accumulation parameters such as drift, threshold, and non-decision time, while the timeout process controls the temporal probability that the trial remains response-eligible.

## Base response process

Let

[
T_R
]

denote the first-passage time of the response process. For a one-choice Wald model,

[
T_R \sim \operatorname{Wald}(a, v, \sigma),
]

where (a) is the response threshold, (v) is drift rate, and (\sigma) is the diffusion coefficient. Observed reaction time is

[
RT = t_0 + T_R,
]

where (t_0) is non-decision time.

The standard model predicts an observed response whenever (T_R) occurs before the experimental cutoff (U - t_0). This can create difficulty when omission rates are non-zero but observed reaction times show little or no density near the cutoff.

## Independent timeout process

Introduce an independent timeout process

[
T_K,
]

where (K) denotes killing, expiry, or abandonment of the response opportunity. A response is observed only if

[
T_R < T_K.
]

If (T_K < T_R), the trial produces an omission or, in a choice model, a time-based guess.

For a detection task with upper response window (u = U - t_0), the probability of observing a hit is

[
P(\text{hit before } u)
= \int_0^u f_R(t) S_K(t),dt,
]

where (f_R(t)) is the response-time density of the evidence process and

[
S_K(t)=P(T_K>t)
]

is the survival function of the timeout process.

The omission probability is therefore

[
P(\text{omission by } u)
= 1 - \int_0^u f_R(t) S_K(t),dt.
]

## Exponential kill process

The simplest timeout model assumes

[
T_K \sim \operatorname{Exponential}(\lambda).
]

Then

[
S_K(t)=e^{-\lambda t},
]

and the observed hit density is

[
f_{\text{hit}}(t)
= f_R(t)e^{-\lambda t}.
]

The parameter (\lambda) is the constant hazard of response-opportunity expiry:

[
h_K(t)=\lambda.
]

This process selectively downweights long response times. It therefore converts some otherwise late responses into omissions while preserving the shape of the early and modal response-time distribution.

For a Wald response process, the eventual hit probability is available through the Laplace transform of the Wald distribution:

[
L(s)=\mathbb{E}\left[e^{-sT_R}\right]
= \exp\left[
\frac{a}{\sigma^2}
\left(
v - \sqrt{v^2 + 2\sigma^2 s}
\right)
\right].
]

Thus, under exponential killing,

[
P(\text{eventual hit})=L(\lambda),
]

and

[
P(\text{eventual omission})=1-L(\lambda).
]

With a finite cutoff (u), define

[
v_\lambda = \sqrt{v^2 + 2\sigma^2\lambda}.
]

Then

[
f_W(t\mid a,v,\sigma)e^{-\lambda t}
===================================

\exp\left[
\frac{a}{\sigma^2}(v-v_\lambda)
\right]
f_W(t\mid a,v_\lambda,\sigma),
]

so

[
P(\text{hit before }u)
======================

\exp\left[
\frac{a}{\sigma^2}(v-v_\lambda)
\right]
F_W(u\mid a,v_\lambda,\sigma).
]

This provides a fully analytic censored likelihood for the killed Wald model.

## Erlang-2 kill process

A limitation of the exponential kill process is its constant hazard. This implies that expiry can occur immediately after stimulus onset with the same instantaneous risk as later in the trial.

A minimal extension uses a two-stage Erlang timeout process:

[
T_K \sim \operatorname{Erlang}(2,\lambda).
]

Its survival function is

[
S_K(t)=e^{-\lambda t}(1+\lambda t),
]

and its density is

[
f_K(t)=\lambda^2 t e^{-\lambda t}.
]

The corresponding hazard is

[
h_K(t)=\frac{\lambda^2 t}{1+\lambda t}.
]

This hazard satisfies

[
h_K(0)=0,
]

and increases monotonically toward (\lambda). The process therefore implements a simple form of increasing time pressure while preserving a compact analytic structure.

For a detection task, the observed hit density becomes

[
f_{\text{hit}}(t)
= f_R(t)e^{-\lambda t}(1+\lambda t).
]

The eventual hit probability is

[
P(\text{eventual hit})
= \mathbb{E}\left[e^{-\lambda T_R}(1+\lambda T_R)\right].
]

Using the Laplace transform (L(s)), this is

[
P(\text{eventual hit})
= L(\lambda)-\lambda L'(\lambda).
]

For the Wald model,

[
L'(\lambda)
= -\frac{a}{\sqrt{v^2+2\sigma^2\lambda}}L(\lambda),
]

so

[
P(\text{eventual hit})
======================

L(\lambda)
\left[
1 + \frac{\lambda a}{\sqrt{v^2+2\sigma^2\lambda}}
\right].
]

The finite-cutoff probability is

[
P(\text{hit before }u)
= \int_0^u f_W(t\mid a,v,\sigma)e^{-\lambda t}(1+\lambda t),dt.
]

Because

[
f_W(t\mid a,v,\sigma)e^{-\lambda t}
===================================

C_\lambda f_W(t\mid a,v_\lambda,\sigma),
]

where

[
C_\lambda=
\exp\left[
\frac{a}{\sigma^2}(v-v_\lambda)
\right],
]

the finite-cutoff probability can be written as

[
P(\text{hit before }u)
= C_\lambda
\left[
F_W(u\mid a,v_\lambda,\sigma)

* \lambda M_W(u\mid a,v_\lambda,\sigma)
  \right],
  ]

where

[
M_W(u\mid a,v_\lambda,\sigma)
= \int_0^u t f_W(t\mid a,v_\lambda,\sigma),dt
]

is the lower incomplete first moment of the transformed Wald distribution.

Let

[
q=v_\lambda,
\qquad
\mu_q=\frac{a}{q},
\qquad
\eta=\frac{a^2}{\sigma^2}.
]

The Wald CDF is

[
F_W(u\mid a,q,\sigma)
=====================

\Phi\left(\frac{q u-a}{\sigma\sqrt{u}}\right)
+
\exp\left(\frac{2aq}{\sigma^2}\right)
\Phi\left(-\frac{q u+a}{\sigma\sqrt{u}}\right),
]

and the lower incomplete first moment is

[
M_W(u\mid a,q,\sigma)
=====================

\mu_q
\left[
\Phi\left(\frac{q u-a}{\sigma\sqrt{u}}\right)
---------------------------------------------

\exp\left(\frac{2aq}{\sigma^2}\right)
\Phi\left(-\frac{q u+a}{\sigma\sqrt{u}}\right)
\right].
]

Therefore the Erlang-2 killed Wald has a closed-form finite-window hit probability:

[
P(\text{hit before }u)
= C_\lambda
\left{
F_W(u\mid a,q,\sigma)

* \lambda M_W(u\mid a,q,\sigma)
  \right},
  \qquad q=v_\lambda.
  ]

## Choice-response extension with time-based guessing

For a choice model with evidence-based response densities (g_i(t)) and total evidence survival

[
S_E(t)=P(T_R>t),
]

the timeout process may produce either omissions or guesses.

If a timeout produces a guess with response probabilities (\pi_i), then the response density for option (i) is

[
p_i(t)
======

g_i(t)S_K(t)

* \pi_i f_K(t)S_E(t).
  ]

For exponential killing,

[
p_i(t)
======

g_i(t)e^{-\lambda t}

* \pi_i \lambda e^{-\lambda t}S_E(t).
  ]

For Erlang-2 killing,

[
p_i(t)
======

g_i(t)e^{-\lambda t}(1+\lambda t)

* \pi_i \lambda^2 t e^{-\lambda t}S_E(t).
  ]

If the timeout produces omissions rather than guesses, the second term is omitted from the response density and contributes instead to the omission probability.

With an upper response window (u), the no-response probability is

[
P(\text{no response by }u)
==========================

S_E(u)S_K(u)
+
\int_0^u f_K(t)S_E(t),dt
]

when timeouts produce omissions. If timeouts always produce guesses, then the no-response probability reduces to

[
P(\text{no response by }u)=S_E(u)S_K(u),
]

because either the evidence process or the timeout process would have generated an observable response before (u).

## Interpretation

The evidence accumulation process determines the ordinary response-time distribution and retains its standard psychological interpretation. The timeout process determines how long the trial remains response-eligible.

The exponential model assumes a constant risk of expiry:

[
h_K(t)=\lambda.
]

The Erlang-2 model assumes expiry risk starts at zero and increases over time:

[
h_K(t)=\frac{\lambda^2 t}{1+\lambda t}.
]

The Erlang-2 process is therefore a minimal analytic approximation to a timing accumulator. It introduces increasing temporal pressure without requiring a full additional accumulation process with its own threshold, drift, and diffusion parameters.

In detection tasks, the timeout process provides a principled mechanism for omissions without requiring a long unobserved response-time tail. In choice tasks, the same mechanism can be used to generate time-based guesses, omissions, or both.
